\title{Parameter-Efficient Orthogonal Finetuning via Factorization}

\begin{abstract}
The widespread adoption of large foundation models is hindered by the substantial computational cost associated with training them from the ground up. Therefore, developing approaches that can effectively and efficiently tailor these large models to new tasks has become essential. This work focuses on Orthogonal Finetuning (OFT), a theoretically motivated strategy for adapting models to downstream applications. Although OFT offers promising generalization, it is hampered by the significant number of trainable parameters resulting from the dimensionality of orthogonal matrices. To mitigate this limitation, we re-examine OFT through the perspective of information transfer, leading to the formulation of several criteria for achieving improved parameter-efficiency. Drawing inspiration from the Cooley-Tukey fast Fourier transform’s capacity for efficient information flow, we introduce a compact orthogonal parameterization using butterfly structures. By integrating this structure into OFT, we present Orthogonal Butterfly~(BOFT), a new finetuning technique that achieves parameter-efficiency while generalizing OFT as a special instance. We evaluate the effectiveness of BOFT through comprehensive empirical experiments, demonstrating its utility in adapting large-scale vision transformers, language models, and text-to-image diffusion models to an array of vision and language tasks.
\end{abstract}

\section{Introduction}

Models such as ChatGPT~\cite{brown2020language,chen2021evaluating} and Stable Diffusion~\cite{rombach2022high} exemplify the outstanding generalization achieved by large foundation models. This leap in model capabilities has been accompanied by a rapid increase in the scale of these models

(\emph{e.g.}, GPT-3 contains approximately 175B parameters), making it ever more impractical to perform training from scratch. Consequently, for continued progress in this area, it is crucial to develop techniques that enable the adaptation of existing foundation models to specific downstream tasks without the necessity of retraining. To this end, the research community has explored three major adaptation paradigms: \emph{model finetuning}~\cite{hulora2022,zaken2022bitfit,guo2020parameter,raffel2020exploring,qiu2023controlling,zhang2022adaptive,chavan2023one}, which keeps the model structure intact and tunes only selected parameters; \emph{adapter tuning}~\cite{rebuffi2017learning,houlsby2019parameter,pfeiffer2020adapterfusion,lin2020exploring,he2021towards}, which augments the model with extra trainable modules; and \emph{prompt tuning}~\cite{li2021prefix,lester2021power}, which leverages additional trainable prefix tokens as part of the input sequence. Among these strategies, model finetuning is particularly appealing due to its straightforward implementation and the absence of additional inference overhead.

Model finetuning fundamentally aims to maintain a close resemblance between the pretrained and finetuned models, as measured by specific criteria, to retain the knowledge obtained during pretraining. Commonly, finetuning strategies employ low learning rates, a pragmatic choice that limits deviation from the pretrained state. Given the inherently structured arrangement of weight matrices, an alternative and more systematic method is to conserve the relational properties—such as the angles between neuron vectors—within these matrices. This perspective underpins the development of Orthogonal Finetuning~(OFT)~\cite{qiu2023controlling}, a framework wherein all neurons in a given layer are subject to transformation by an identical orthogonal matrix, thereby guaranteeing that the pairwise angles between neurons are exactly maintained throughout finetuning. OFT has shown favorable generalization and efficient convergence in the finetuning of text-to-image diffusion models~\cite{qiu2023controlling}. Nonetheless, the use of full orthogonal matrices can entail a substantial number of trainable parameters because of their high dimensionality. To mitigate this, OFT adopts a block-diagonal design to curtail parameter count. Yet, this greater efficiency is not without cost—block-diagonal orthogonal matrices enforce a predetermined sparsity structure and each block operates independently. While this configuration conserves parameters, it may lead to less flexible representations (\emph{e.g.}, decreasing the capacity of the transformation to mimic traditional linear mappings) and, importantly, the optimal choice of sparsity pattern for such matrices remains unresolved.

To tackle this challenge, it is crucial to produce a dense orthogonal matrix while maintaining a high degree of parameter efficiency. Although constructing a dense orthogonal matrix of size $d$ generally necessitates $\mathcal{O}(d^2)$ parameters, our strategy departs from conventional approaches by assembling such a matrix through the product of multiple factorized sparse matrices. This methodology is driven by the realization that parameter count can be minimized by sacrificing computational speed for reduced memory requirements. In our method, since the orthogonal matrix is composed as a sequence of sparse matrix multiplications, these operations must be executed iteratively during training. To contextualize our matrix decomposition, we recast the construction of a dense orthogonal matrix as a task of information propagation. In this paradigm, the multiplication of sparse matrices to yield a dense orthogonal matrix can be seen as the transfer of information across the nodes of a grid-structured graph. The adoption of this information transmission model provides flexibility in developing a variety of sparse matrix factorization schemes, balancing parameter budget and the ability to represent dense orthogonal matrices effectively.

Our framework for efficient parameterization leverages the butterfly structures integral to the Cooley-Tukey fast Fourier transform algorithm~\cite{cooley1965algorithm}, where communication between graph nodes is achieved in a highly efficient manner~\cite{klasing1994broadcasting}. The butterfly graph topology, as employed in Cooley-Tukey, directly motivates what is known as butterfly factorization for sparse matrix decomposition. With this factorization scheme, it is possible to synthesize a $d$-dimensional dense matrix as the product of only $\mathcal{O}(\log d)$ sparse matrices, with each containing $\mathcal{O}(d)$ nonzero elements. Enforcing the orthogonality of each constituent sparse matrix leads to a representation of the full orthogonal matrix with just $\mathcal{O}(d\log d)$ parameters—an appreciable reduction from the naïve quadratic scaling. Capitalizing on this structure, we introduce Orthogonal Butterfly (BOFT), a novel finetuning method that is highly parameter-efficient. BOFT generalizes OFT, accommodating it as a specific instance within a broader orthogonal finetuning scheme. Both block-diagonal and butterfly matrix structures share an inherent sparsity, which is the principal mechanism by which they reduce the parameter count required for orthogonal matrices. Notably, this stands in contrast to the low-rank approach taken by LoRA~\cite{hulora2022}; while LoRA and our method both rely on structured matrices for parameter savings, they stem from two of the primary categories—low-rank and sparse matrices—identified as efficient matrix structures~\cite{candes2011robust}.

In contrast to OFT, which leverages a block-diagonal parameterization to balance expressivity and regularity, BOFT adopts the butterfly architecture to enable a more gradual and continuous transition between elements of the full orthogonal group (encompassing expressivity) and the identity matrix (reflecting regularity). This approach allows us to identify a more compact subset within the orthogonal group that suffices for downstream applications. Given the extensive application of butterfly structures in a variety of efficient linear transformations, including the discrete Fourier and discrete cosine transforms, we posit that incorporating such structure imparts an implicit inductive bias, which facilitates improved generalization and serves as a safeguard against overfitting. The rationale behind this claim is that butterfly parameterizations naturally represent many classical linear operators, an expressivity the block-diagonal design in OFT cannot match. Our main contributions are as follows:

\begin{itemize}[leftmargin=*,nosep]
    \item We tackle the issue of parameter efficiency in orthogonal finetuning through an innovative information transmission framework. This paradigm reframes the problem of designing dense, parameter-efficient orthogonal matrices as an information flow challenge on a graph with grid-like connectivity.
    \item Building on the butterfly factorization techniques employed in the Cooley-Tukey algorithm, we introduce Orthogonal Butterfly, a new approach for parameter-efficient orthogonal finetuning within our information transmission model.
    \item We deliver several theoretical justifications for BOFT's ability to substantially compress the parameter space without sacrificing representational capacity or stability during training. The factorized form of the matrices in BOFT also enables a distinctive method for interpolating weights (refer to Figure~\ref{fig:boft_interpolation}).
    \item We pioneer the application of orthogonal finetuning to a broad array of tasks that extend beyond controllable text-to-image generation~\cite{qiu2023controlling}, establishing BOFT as a highly versatile tool for model adaptation. To demonstrate its wide applicability, we apply BOFT across diverse domains such as computer vision and natural language processing, where it consistently surpasses state-of-the-art alternatives in both parameter efficiency and generalization performance.
\end{itemize}

\section{Related Work}

\textbf{Parameter-efficient finetuning (PEFT)}. With the growing scale and capability of foundation models (\emph{e.g.}, \cite{brown2020language,radford2021learning,rombach2022high,kirillov2023segment}), there has been significant attention paid to efficient methods for adapting these models to downstream tasks while modifying as few parameters as possible~\cite{treviso2023efficient,ding2023parameter,lialin2023scaling}. A range of PEFT strategies has been proposed~\cite{houlsby2019parameter,aghajanyan2020intrinsic,hulora2022,edalati2022krona,wang2022adamix,gheini2021cross,zaken2022bitfit,guo2020parameter,sung2021training,ansell2022composable,lester2021power,li2021prefix,vu2022spot,he2021towards,mao2021unipelt,karimi2021compacter,liu2022few,sung2022lst,chen2022parameter,jia2022visual,chen2022adaptformer,zhang2022neural,jie2023fact,lian2022scaling,luo2023towards}, of which reparameterization-based techniques~\cite{aghajanyan2020intrinsic,hulora2022,edalati2022krona} are particularly pertinent to this study. In particular, LoRA~\cite{hulora2022} enhances a pretrained weight matrix by incorporating the product of two low-rank matrices, delivering competitive outcomes for natural language processing applications. While LoRA utilizes a fixed rank for the added matrices, more recent methods \cite{zhang2022adaptive,valipour2022dylora,zhang2023increlora} adaptively select layer-specific ranks to more effectively distribute the available parameters. Owing to their simplicity, low-rank reparameterization techniques are widely adopted~\cite{dettmers2023qlora,zi2023delta,chavan2023one}. Drawing motivation from the use of hyperspherical energy to explain generalization~\cite{liu2018learning,liu2021orthogonal}, orthogonal finetuning was introduced in \cite{qiu2023controlling} as an alternative and effective adaptation mechanism for text-to-image diffusion models. This approach optimizes an orthogonal matrix to rotate the neurons within the same layer, resulting in improved generalization and steadier training relative to LoRA. However, OFT generally entails a greater number of trainable parameters than LoRA, highlighting the value of pursuing higher parameter efficiency in OFT. Furthermore, it remains to be established whether OFT is extensible to broader adaptation scenarios beyond text-to-image diffusion models~\cite{qiu2023controlling}. The BOFT method addresses the parameter efficiency limitations of OFT by employing butterfly factorization, enabling us to validate the broader applicability and effectiveness of orthogonal finetuning in general adaptation settings.

\textbf{Butterfly structures}. The radix-2 Cooley-Tukey algorithm reduces the $N$-point discrete Fourier transform problem into two subproblems of size $\frac{N}{2}$ recursively, resulting in the emergence of butterfly patterns, which can be formalized as a product of multiple sparse matrices—collectively termed a butterfly matrix. Such butterfly matrices have served as tools for parameterizing orthogonal matrices (circumventing the need for pivoting in Gaussian elimination and enhancing computational efficiency~\cite{parker1995random}), stabilizing the learning of recurrent neural networks~\cite{jing2017tunable}, and improving kernel approximations~\cite{munkhoeva2018quadrature}. Research by \cite{dao2019learning,dao2019kaleidoscope} has utilized butterfly parameterizations to learn rapid linear transforms, while \cite{chen2022pixelated,dao2022monarch} leverage butterfly matrices to facilitate efficient neural network training. Applications for butterfly structures further include fast matrix-vector operations~\cite{michielssen1996multilevel,o2010algorithm}, compressing matrix representations~\cite{li2015butterfly}, and optimizing network communication~\cite{klasing1994broadcasting,li2003linear,rajasingh2016transmission}. Unlike earlier efforts, this work seeks to leverage butterfly structures specifically to boost the parameter efficiency of OFT.

\section{An Information Transmission View on Orthogonal Finetuning}

We begin by outlining foundational concepts of OFT. In order to adapt the pretrained weight matrix, OFT expresses the updated weights as a product of a trainable orthogonal matrix with the static pretrained weights. Unlike LoRA—which implements adaptation through the addition of a low-rank matrix—OFT leverages a multiplicative orthogonal transformation to refine the weights. LoRA achieves parameter efficiency via its low-rank formulation; by comparison, the original OFT adopts a block-diagonal orthogonal parameterization~\cite{qiu2023controlling}, where smaller block sizes increase parameter efficiency. Figure~\ref{fig:comp_lora} illustrates this comparison. The underlying motivation for using orthogonal transformations during finetuning is to retain pair-wise neuron angles~\cite{liu2018learning,liu2021orthogonal,qiu2023controlling}, thus maintaining a large portion of the semantic information encoded during pretraining.

More specifically, OFT introduces an orthogonal matrix $\bm{R}\in\mathbb{R}^{d\times d}$ for a pretrained linear layer $\bm{W}^0\in\mathbb{R}^{d\times n}$, updating the computation from $\bm{z}=(\bm{W}^0)^\top\bm{x}$ to $\bm{z}=(\bm{R}\bm{W}^0)^\top\bm{x}$, with $\bm{x}\in\mathbb{R}^{d}$ and $\bm{z}\in\mathbb{R}^{n}$ representing the input and output, respectively. For zero initialization, OFT sets $\bm{R}$ initially as the identity matrix. To ensure $\bm{R}$ maintains its orthogonality throughout finetuning, we utilize the Cayley transform following~\cite{qiu2023controlling,liu2021orthogonal}, specifically $\bm{R}=(\bm{I}+\bm{Q})(\bm{I}-\bm{Q})^{-1}$, where $\bm{Q}$ is a skew-symmetric matrix, i.e., $\bm{Q}=-\bm{Q}^\top$. To further promote parameter efficiency, $\bm{R}$ is structured as block-diagonal: $\text{diag}(\bm{R}_1,\bm{R}_2,\cdots,\bm{R}_r)$, with each $\bm{R}_i\in\mathbb{R}^{b\times b}$ being a small orthogonal block and $br=d$. However, this block-diagonal approach introduces the implicit assumption that the neuron dimension (i.e., a column of $\bm{W}^0$) is partitioned into $r$ segments, each independently transformed by distinct orthogonal matrices. Even though the block-diagonal method has been empirically validated, dividing a neuron’s dimensions based solely on their indices lacks principled justification, highlighting the potential of using dense orthogonal matrices. This observation prompts an important question: \emph{Is it possible to devise a dense orthogonal matrix that preserves parameter efficiency?}

To tackle this issue, we introduce the idea of constructing a dense orthogonal matrix by composing several sparse orthogonal matrices. To systematically investigate the structure of sparsity in orthogonal matrix factorizations, we reinterpret the task of synthesizing a dense orthogonal matrix in OFT as an instance of information transmission. More concretely, forming a dense matrix $\bm{R}\in\mathbb{R}^{d\times d}$ via the sequential multiplication of $m$ square matrices as $\thickmuskip=2mu \medmuskip=2mu \bm{R}=\bm{B}_m\bm{B}_{m-1}\cdots\bm{B}_1$ can be analogized to relaying information through a grid comprised of $d\times (m+1)$ nodes, as depicted in Figure~\ref{fig:info_trans}. The rationale for this information flow analogy stems from viewing a dense $d\times d$ matrix as a fully connected mapping from one set of $d$ nodes to another. Specifically, for the matrix $\bm{R}$, a zero entry $R_{ij}$ signals that it is impossible to propagate information from the $j$-th to the $i$-th node, whereas a nonzero entry permits such transmission. Accordingly, decomposing $\bm{R}$ into the product $\bm{B}_m\bm{B}_{m-1}\cdots\bm{B}_1$ equates to a process of stepwise information exchange dictated by the connectivity defined by each $\bm{B}_i$. Information thus propagates initially according to $\bm{B}_1$ and ultimately via $\bm{B}_n$. As an explicit illustration in Figure~\ref{fig:info_trans}, consider the factorization $\bm{R}=\bm{B}_5\bm{B}_4\bm{B}_3\bm{B}_2\bm{B}_1$, where both the sparsity patterns and the associated connectivity graph are visualized. The shown graph in Figure~\ref{fig:info_trans} corresponds to the expansion of this matrix product. Within this induced graph, each $\bm{B}_i$ operates as the connectivity map between the $i$-th and $(i+1)$-th sets of nodes; the $(j_1,j_2)$ entry in $\bm{B}_i$ indicates the presence or absence of a directed link from node $j_2$ at level $i$ to node $j_1$ at level $(i+1)$ (where zero indicates no link). In order for $\bm{B}_5\bm{B}_4\bm{B}_3\bm{B}_2\bm{B}_1$ to yield a dense matrix, each node in the initial layer must be able to transmit information to every node in the sixth layer. By limiting our attention to the factorization $\bm{R}=\bm{B}_4\bm{B}_3\bm{B}_2\bm{B}_1$, which relates to source nodes at the first level and destination nodes at the fifth, it becomes apparent, for instance, that node 1 does not have a transmission path to node 3. Consequently, the pattern of zeros in $\bm{B}_4\bm{B}_3\bm{B}_2\bm{B}_1$ features a zero entry at $(3,1)$. Similar reasoning applies for $\bm{R}=\bm{B}_3\bm{B}_2\bm{B}_1$ and $\bm{R}=\bm{B}_2\bm{B}_1$, establishing a general correspondence between the connectivity of the induced graph and the sparsity layout. More broadly, achieving density in a matrix $\bm{R}\in\mathbb{R}^{d\times d}$ through factorization as $\bm{B}_m,\ldots,\bm{B}_1$ necessitates that these matrices correspond to a collection of directed edges overlaying a $d\times (m+1)$ grid, where edges can only connect pairs of adjacent levels (that is, adjacent columns), such that every node from the first layer is capable of reaching every node in the $(m+1)$-th layer.

Adopting the information transmission perspective allows for deeper insight into the sparsity structure characteristic of factorization matrices within OFT. In Figure~\ref{fig:blk_diag}, the block-diagonal arrangement of $\bm{R}$ from the canonical OFT is depicted. While this pattern succeeds in lowering the parameter count, it fails to yield a fully dense matrix $\bm{R}$. Our objective is to assemble a dense orthogonal matrix from $m$ sparse orthogonal matrices while minimizing the number of trainable parameters. The information transmission paradigm leads us to two primary requirements: \emph{(i)} \textbf{dense connectivity}, which ensures that every node at the initial layer is linked, via some path, to all nodes in the terminal layer; and \emph{(ii)} \textbf{minimum free edges}, where, given the orthogonality requirement, we seek the smallest possible number of total edges. It is important to recognize that orthogonality imposes specific restrictions on the allowable connections between neighboring layers. Specifically, to ensure each matrix $\bm{B}_i$ retains full rank—a prerequisite for orthogonality—there must be $d$ connections that form a bijection between the nodes of the $i$-th and $(i+1)$-th layers, dictating at least $d$ edges between each pair of adjacent layers (as illustrated, for instance, with $d=4$ in Figure~\ref{fig:info_trans}). These $d$ mandatory connections, essential for preserving orthogonality, are not counted among the edges contributing trainability, as the corresponding entries in the matrix are not learnable (for example, the non-zero entries in a $d \times d$ orthogonal matrix with $d$ non-zeros can only take the values $\pm 1$). Due to the orthogonality condition, which reduces the required parameters compared to unconstrained matrices, there is increased dependency among edges. For instance, in a $2\times 2$ orthogonal matrix, a zero entry at position $(1,1)$ compels a zero at $(2,2)$ as well (i.e., the absence of one connection implies the absence of another). Consequently, the orthogonality property can necessitate adding or removing specific edges within a given configuration. Through examining the nonzero entry pattern in the resulting orthogonal matrix, the information transmission model becomes especially instructive in the context of OFT, since our concern is restricted to the placement of trainable, nonzero elements in $\bm{R}$ rather than their exact numerical values.

A straightforward approach to establishing dense connectivity between two layers requires $\mathcal{O}(d^2)$ connections (\emph{i.e.}, a full dense orthogonal matrix), amounting to $d^2-d$ total edges (e.g., for $d=4$, this equals 12 edges). As shown in Figure~\ref{fig:info_trans}, one can construct a viable matrix factorization with only $10$ edges, fewer than those in a dense orthogonal matrix. This graphical framework permits an analysis of OFT’s parameter efficiency, facilitating the identification of valid factorizations. We draw on the distinctive architecture of butterfly graphs~\cite{cooley1965algorithm}, well-known from the Cooley-Tukey algorithm, which efficiently establishes dense connectivity between $d$ source and $d$ target nodes with merely $\mathcal{O}(d\log d)$ edges. For instance, Figure~\ref{fig:info_trans} employs $10$ edges, while the butterfly network requires only $8$ to accomplish full connectivity. In what follows, we discuss how the butterfly topology enhances parameter efficiency.

\section{Orthogonal Parameterization by Butterfly Factorization}

The butterfly pattern, originating in the Cooley-Tukey algorithm, is integral to executing the fast Fourier transform (FFT). Within the Fourier transform, localized modifications in frequency can lead to widespread alterations in the spatial domain, paralleling the objective of efficient information transfer—enabling every node in the initial layer to communicate with all nodes in the terminal layer. The butterfly network has also gained traction in computer networking~\cite{solihin2015fundamentals,li2011linear} for facilitating swift information transmission. Suppose $k \geq 2$ is a power of $2$; we introduce the \emph{butterfly factor} $\bm{B}^F(k)$ as

\begin{equation}\label{eq:but_factor}
\begin{aligned}
    \bm{B}^F(k)=\begin{bmatrix}
    \text{diag}(\bm{d}_1) & \text{diag}(\bm{d}_2) \\
    \text{diag}(\bm{d}_3) & \text{diag}(\bm{d}_4)
    \end{bmatrix}\in\mathbb{R}^{k\times k},
\end{aligned}
\end{equation}

where each $\bm{d}_i\in\mathbb{R}^{\frac{k}{2}}$ for $i=1,\ldots,4$ represents a vector. Given $d=2^N$, the $d$-dimensional \emph{butterfly matrix} $\bm{B}(d)\in\mathbb{R}^{d\times d}$ is defined recursively:

\begin{equation}
\begin{aligned}
    \!\!\bm{B}(d)=\tilde{\bm{B}}(d,d)\!\cdot\!\begin{bmatrix}
    \bm{B}_1(\frac{d}{2})\!\!\!\!\! & \bm{0} \\
    \bm{0}\!\!\!\!\! &\bm{B}_2(\frac{d}{2})
    \end{bmatrix}= \tilde{\bm{B}}(d,d)\tilde{\bm{B}}(d,\frac{d}{2})\cdots\tilde{\bm{B}}(d,2),
\end{aligned}
\end{equation}

Here, $\bm{B}_1(\frac{d}{2})$ and $\bm{B}_2(\frac{d}{2})$ represent two butterfly matrices, each with dimensionality $\frac{d}{2}$. We introduce the \emph{butterfly component} as $\thickmuskip=2mu \medmuskip=2mu \tilde{\bm{B}}(d,k)=\text{diag}(\bm{B}^F_1(k),\cdots,\bm{B}^F_{d/k}(k))$, which is a $d\times d$ block-diagonal matrix comprising blocks of size $k$. In this structure, each $\bm{B}^F_i(k)$ denotes a butterfly factor as specified in Equation~\ref{eq:but_factor}. With these foundations, we can now employ the butterfly matrix as a means of representing orthogonal matrices. Ensuring the orthogonality of every multiplicative factor $\tilde{\bm{B}}(d,k)$ within the overall butterfly matrix $\bm{B}(d)$ is the central requirement. To illustrate, we examine the block-diagonal form $\tilde{\bm{B}}(d,2)$, which has $2\times2$ blocks. Each of these blocks can be parameterized as orthogonal by applying the Cayley transform (or a two-dimensional rotation), consistent with the approaches in \cite{liu2021orthogonal,qiu2023controlling}. The sparsity structure (non-zero pattern) of any butterfly component can be interpreted as a permutation of the sparsity in $\tilde{\bm{B}}(d,2)$, enabling all butterfly components to be parametrized as orthogonal in a straightforward manner. Consequently, this leads to a highly efficient parameterization of orthogonal matrices via collections of $2\times 2$ orthogonal blocks. Building on the methodology of \cite{chen2022pixelated}, we extend butterfly matrices to define the block butterfly component $\tilde{\bm{B}}^b(d,k)$, where now every entry of $\bm{d}_i$ is a $b\times b$ matrix. Orthogonality of the block butterfly component $\tilde{\bm{B}}^b(d,2)$ is ensured by parameterizing each $2b\times 2b$ block as an orthogonal matrix. For other butterfly components $\tilde{\bm{B}}^b(d,k)$ with $k>2$, the non-zero pattern is realized as a block-wise permutation of that in $\tilde{\bm{B}}^b(d,2)$, so the orthogonality condition can be satisfied similarly. Synthesizing all these aspects, the BOFT forward pass is given by

\begin{equation}
    \begin{aligned}\nonumber
    \bm{z}=\big(\bm{R}(m,b)\cdot\bm{W}^0\big)^\top\bm{x},~~~~\text{s.t.}~\bigg{\{}\bm{R}(m,b)=\prod_{i=1}^m \tilde{\bm{B}}^b_{(d,i)}~~\&~~\underbrace{\big(\tilde{\bm{B}}^b_{(d,j)}\big)^\top\tilde{\bm{B}}^b_{(d,j)}=\tilde{\bm{B}}^b_{(d,j)}\big(\tilde{\bm{B}}^b_{(d,j)}\big)^\top\!=\bm{I}_d}_{\forall j\in [1, m]}\bigg{\}},
    \end{aligned}
\end{equation}

In this setting, for notational clarity, we let $\tilde{\bm{B}}^b(d,2^{m - i+1})$ be expressed as $\tilde{\bm{B}}^b_{(d,i)}$, and let $\bm{I}_d$ denote the $d \times d$ identity matrix. The orthogonal matrix $\bm{R}(m,b) \in \mathbb{R}^{d \times d}$ arises from a sequence of multiplications of orthogonal butterfly components. For brevity, we refer to BOFT equipped with $\bm{R}(m,\frac{b}{2})$ as BOFT($m$, $b$), with the stipulation that $b \geq 2$. If $\thickmuskip=2mu \medmuskip=2mu m=1$, then BOFT($1$, $b$) coincides with the block-diagonal OFT~\cite{qiu2023controlling} of block size $b$. Moreover, setting $b=d$ in BOFT($1$, $d$) yields the classic OFT~\cite{qiu2023controlling}, which employs a full, unconstrained orthogonal matrix. In the configuration BOFT($\log \frac{2d}{b}$, $b$), $\bm{R}$ is instantiated as the block butterfly matrix $\bm{B}^{\frac{b}{2}}(d)$, giving rise to a dense orthogonal matrix. In a more general case, BOFT($m$, $b$) features $\frac{1}{2}(b-1)dm$ learnable parameters when adapting a linear transformation of dimension $d\times n$. Employing the butterfly structure ($m = \log d, b=2$), BOFT requires only $\mathcal{O}(d\log d)$ parameters. In comparison, conventional OFT that uses a full dense orthogonal matrix needs $\mathcal{O}(d^2)$ parameters, whereas the block-diagonal OFT with $r$ blocks entails $\mathcal{O}(bd)$ parameters. Thus, generating a dense orthogonal matrix via standard OFT necessitates the block size $b=d$, while BOFT can realize a dense orthogonal transformation for arbitrary $b$.

\textbf{Identity Initialization in BOFT}. Typical fine-tuning protocols commence with weights matching the original pretrained model, minimizing deviation during adaptation. For instance, LoRA applies zero initialization to the low-rank parameters. Analogously, in BOFT, every butterfly component is initialized as the identity matrix (that is, initializing the underlying skew-symmetric matrix in the Cayley transform to zeros).

\textbf{Multiplicative Dropout Mechanism in BOFT}. LoRA~\cite{hulora2022} incorporates a Dropout mechanism to counteract overfitting by regularizing the low-rank updates. Standard Dropout~\cite{srivastava2014dropout} aligns well with the LoRA structure but is unsuitable for BOFT because BOFT's parameter updates are multiplicative rather than additive. We address this limitation by proposing a multiplicative Dropout approach tailored for BOFT. Since the orthogonal matrix $\bm{R}(m,b)$ is assembled from $m$ butterfly components, which themselves can be rearranged into $2b \times 2b$ block-diagonal orthogonal structures, our dropout process operates as follows: a proportion $p_1$ of the butterfly components is randomly chosen, and within each component, $p_2$ of the diagonal blocks are selected. These selected blocks are then substituted by identity matrices.

\section{Intriguing Insights and Discussions}

\textbf{Expressivity of BOFT}. The butterfly architecture, in combination with appropriate permutations, is capable of exactly representing numerous traditional fast linear transforms~\cite{dao2019learning,dao2019kaleidoscope} (\emph{e.g.}, fast Fourier transform, Hadamard transform). However, the degree to which our orthogonal butterfly matrix can approximate an arbitrary orthogonal matrix is still not well characterized. To address this, we conducted a simulation to fit a randomly generated dense orthogonal matrix~\cite{anderson1987generation} of dimension $512\times 512$. Figure~\ref{fig:para_eff} presents the results averaged over 10 separate random seeds. Here, the y-axis measures the approximation error, while the x-axis represents the count of effective trainable parameters. Each curve of a given color corresponds to BOFT with a specific block size; notably, the point at the far left of each curve depicts the error from BOFT($1$,$b$) (\emph{i.e.}, standard block-diagonal OFT with block size $b$). BOFT consistently demonstrates improved parameter efficiency in comparison to OFT. For instance, BOFT($9$,$2$) offers greater expressivity than BOFT($1$,$16$) despite using substantially fewer parameters. Likewise, configurations with smaller $b$ and larger $m$ tend to be more parameter-efficient, as evidenced by BOFT($6$,$4$) matching the approximation error of BOFT($2$,$16$) while employing many fewer parameters. This can be attributed to the butterfly matrix’s imposition of a richer structure within the orthogonal group than the block-diagonal approach, making BOFT provably more expressive than an OFT of equivalent block size.

\begin{theorem}[Expressivity of BOFT]\label{thm:exp_boft}
    For a fixed block size, BOFT exhibits strictly greater expressiveness than OFT. To achieve universality over all orthogonal matrices of size $d$, one may take a product of butterfly matrices as $\bm{B}_{d-1,1}(d)\bm{B}^\top_{d-1,2}(d)\cdots\bm{B}_{1,1}(d)\bm{B}^\top_{1,2}(d)$, with each $\bm{B}_{i,j}(d),\forall i,\forall j$ a butterfly matrix.
\end{theorem}

Based on Theorem~\ref{thm:exp_boft}, one can generalize BOFT by forming the final orthogonal matrix as $\thickmuskip=2mu \medmuskip=2mu \bm{R}^G(m_1,b_1,m_2,b_2,l)=\bm{R}_{l,1}(m_1,b_1)\bm{R}_{l,2}^T(m_2,b_2)\cdots\bm{R}_{1,1}(m_1,b_1)\bm{R}_{1,2}^T(m_2,b_2)$, where $\bm{R}_{i,j}^T(m,b)$ indicates the orthogonal matrices within the BOFT structure. Particularly, setting $m_1=m_2=\log d$, $b_1=b_2=2$, and $l=d-1$, $\bm{R}^G(m_1,b_1,m_2,b_2,l)$ is able to span the full orthogonal group. This hierarchical composition is referred to as a kaleidoscope hierarchy~\cite{dao2019kaleidoscope}. Nevertheless, greater expressivity does not inherently guarantee superior finetuning outcomes—indeed, unrestricted finetuning, though universally expressive, often fails to deliver optimal empirical performance. Achieving a suitable balance between expressiveness and regularization remains central to successful model finetuning. BOFT extends the space of tunable parameters via structured inductive biases, empowering more favorable trade-offs between model flexibility and generalization.

\textbf{Spectral properties}. In terms of spectral characteristics, orthogonal finetuning tends to surpass LoRA, primarily because it maintains the spectral norm of the original pretrained matrix $\bm{W}^0$ exactly. This becomes evident when considering the singular value decomposition: $\bm{W}^0=\bm{U}\bm{\Sigma}\bm{V}^\top$, where $\bm{U}$ and $\bm{V}$ are orthogonal matrices and $\bm{\Sigma}$ contains the singular values along the diagonal. Both OFT and BOFT involve left-multiplying by an orthogonal matrix $\bm{R}$, leading to updated weights of the form $\bm{R}\bm{U}\bm{\Sigma}\bm{V}^\top$. This operation leaves the largest singular value (that is, the spectral norm of $\bm{W}^0$) unchanged. Preservation of this norm has been documented as significantly enhancing both training stability and model generalization~\cite{miyato2018spectral,yoshida2017spectral}. Additional mathematical insights can be found in Appendix~\ref{app:sn_property}.

\textbf{Orthogonal finetuning as learning bilinear similarity}. The BOFT approach can be interpreted through the lens of bilinear similarity learning, where the transformation is expressed as $\bm{w}^0_i\bm{R}\bm{x}$ with $\bm{w}^0_i$ denoting the $i$-th column (or neuron) in $\bm{W}^0$. In this setup, BOFT corresponds to optimizing the bilinear similarity matrix $\bm{R}$ under the constraint that $\bm{R}$ remains orthogonal. This framework has deep ties to distance metric learning~\cite{xing2002distance} as well as classical bilinear forms~\cite{roman2005advanced}.

\textbf{Inductive bias and generalization in BOFT}. In BOFT, the matrix $\bm{R}(m,b)$ often inhabits a particular structured subset within the full orthogonal group, effectively narrowing the hypothesis space. This structural limitation naturally introduces an inductive bias. We posit that the particular inductive bias shaped by butterfly factorization is advantageous for generalization because it embodies structural motifs seen in many standard linear transformations~\cite{dao2019kaleidoscope}, such as the discrete Fourier transform, discrete sine/cosine transform, and the Hadamard transform. Furthermore, the sparse nature of matrix factorization in BOFT is likely to yield additional implicit inductive bias~\cite{gunasekar2017implicit,li2018algorithmic,liu2019neural}.

\textbf{Comparison to butterfly-based sparse training}. Several prior studies~\cite{dao2019learning,dao2019kaleidoscope,chen2022pixelated,dao2022monarch} have explored sparse training methods utilizing butterfly parameterizations. These works generally emphasize reparameterizing weight matrices via the butterfly structure and subsequently training neural networks from the beginning. In contrast, \cite{dao2022monarch} investigates finetuning by first projecting pretrained weights onto a butterfly matrix variant, followed by optimizing the resulting projected elements for specific tasks. BOFT introduces a fundamentally distinct finetuning approach, employing shared weight matrices across layers to transform the model weights.

\input{rebuttal-results}

\section{Concluding Remarks and Limitations}

This paper introduces Orthogonal Butterfly, a broadly applicable and parameter-efficient finetuning approach for foundation models that leverages the butterfly structure. The central principle for enhancing parameter efficiency lies in expressing a dense orthogonal matrix as a product of multiple sparse orthogonal matrices. To facilitate the discovery of practical matrix factorizations, we develop a graphical information transmission framework. Using this perspective, we demonstrate that the butterfly structure is particularly well suited for our goal of sparse orthogonal matrix factorization. Through extensive experiments, we establish the empirical efficacy of BOFT for finetuning large language models, large vision models, and text-to-image generative models, confirming BOFT’s effectiveness as a versatile finetuning technique.

While BOFT yields promising results, certain limitations persist. Due to the construction of the final orthogonal matrix as a product of several orthogonal matrices, the training phase incurs a modest increase in runtime overhead compared to OFT. Strategies to further reduce BOFT’s training cost remain an open area for future work. Nevertheless, once finetuning is complete, the orthogonal matrices learned by BOFT can be integrated directly into the pretrained models, ensuring there is no additional inference overhead. Furthermore, it remains unclear whether the butterfly network represents the most optimal architecture for information transmission. The information transmission framework we present draws inspiration from research in computer networking, where the efficiency of network topologies for information exchange is a major area of investigation. It is plausible that other network structures could lead to even more effective compositions for dense orthogonal matrices.